\documentclass[a4paper, 12pt]{article}
\usepackage{listings}
\usepackage{xcolor}
\usepackage{tikz}

% Pseudocode settings
\lstset{
    language=Python,
    basicstyle=\ttfamily,
    keywordstyle=\color{blue},
    commentstyle=\color{green},
    stringstyle=\color{red},
    frame=single,
    breaklines=true
}

\title{Analise de Dados}
\date{\today}

\begin{document}

\maketitle

\section{Analise Descritiva}

A analise visual de grandes quantidades de dados é uma tarefa bastante complexa. Para auxiliar nessa tarefa normalmente utiliza-se a analise descritiva baseada no tipo de dado que está sendo analisado

\subsection{Univariados}

\begin{itemize}

\item média: muito influenciada por valores acima do normal 

\item moda: para valores categoricos

\item mediana: ordena e retorna o valor central ou a média dos valores centrais  

\item percentil: divide o conjunto em 4 

\item variancia e desvio padrão

\end{itemize}

\subsection{Multivariados}

\begin{itemize}

\item variancia

\item covariancia

\end{itemize}

\section{Pré-Processamento}

\begin{description}

\item[Dados desbalanceados] ara lidar com dados desbalanceados pode-se obter novos dados para a classe minoriária ou adicionar dados artificiais

\item[Limpeza dos dados] Muitas condições podem afetar a qualidade dos dados como, falha humana, limitação dos sensores, ma fé e outros

\item[Dados incompletos] podem ser preenchidos utilizando técnicas estatísicas se a quantidade a ser preenchida não for muito discrepante em relação aos dados que não estão vazios

\item[Dados redundanes] podem confundir o modelo e dar a entender que o dado redundante tem mais importancia em relação aos outros dados, dito isso, é importante tratar esses casos

\item[Ruídos] podem levar ao super-ajuste do modelo mas devem ser caracterizados antes de removidos pois podem ser informações importantes dentro do conjunto

\item[Transformação de dados] Dependendo do método a ser aplicado, valores numéricos devem ser convertidos para categoricos e vice versa

\end{description}

\end{document}
